\section{Extended experimental results}
\label{app:experiments}

This appendix gives the per-family settings, additional comparisons, and measurements summarized in Section~\ref{sec:experiments}.

\subsection{Evaluation protocol}

\paragraph{One acceptance check.} Every returned point, from our solvers and
from every competitor, is accepted only if its relative constraint violation and
its objective gap relative to $1 + |f^\star|$ are both at most the target
tolerance, where $f^\star$ is a high-accuracy reference objective computed once
and outside every timed region. No solver is credited for its own stopping
rule.

\paragraph{A completed field.} Generic solvers (Clarabel, PIQP, HiGHS, OSQP,
SCS, ECOS, ProxQP, PDLP and cuPDLPx) each run in their own interpreter. Where
a structure-specific solver exists it is added: scikit-learn, celer and skglm
for lasso and group lasso, LIBSVM and LIBLINEAR for SVMs, network simplex and
minimum-cost flow for network LPs, and a Frank--Wolfe core-set method for the
minimum enclosing ball.

\paragraph{Timing.} Setup is timed on both sides, imports are excluded on both
sides, baselines run in fresh processes, and repetitions are paired. Head-to-head
speedups (Table~\ref{tab:specialists}) time our method in a fresh process with
structure detection, formulation build, compilation and solve inside one clock
and the persistent compilation cache disabled (the cold first solve), five
repetitions, and also report a second solve in the same process (warm).
Held-out solve rates come from routing runs with a 20\,s budget, where the first
visit includes exploring candidates. Solves run on an RTX~4090 workstation and
on an H100 80\,GB node, and each comparison names its machine.

\subsection{Evolution produces solvers, formulations and operators}

\paragraph{Parameters and schemes: total-variation denoising.}
The first cell is anisotropic TV denoising,
$\min_x \tfrac12\|x - y\|^2 + \lambda\|Dx\|_1$ with $\lambda = 0.1$, on real
clinical retinal optical coherence tomography (OCT) B-scans from the Kermany
dataset \citep{Kermany2018Identifying}: $128\times128$ patches, with the raw
image as $y$, 6 training and 3 held-out patches. A solve terminates at relative
duality gap $\le 10^{-4}$ (Section~\ref{sec:methods}). The baseline is the
fastest of the four base schemes at textbook, typechecked defaults
(Condat--V\~u at $0.346$\,s training SGM-10). Search runs 10 generations of 8
proposals with elitism 1, and champions are evaluated once on the held-out
patches.

\begin{table}[t]
\centering
\small
\begin{tabular}{lllll}
\toprule
config & champion genome & train SGM & held-out SGM & held-out fail \\
\midrule
default (certified) & \texttt{condat\_vu} $\tau{=}0.5,\ \sigma{=}0.186$ & 0.346s & 0.621s & 0/3 \\
random-evolved & \texttt{condat\_vu} $\tau{=}0.024,\ \sigma{=}3.70$ & 0.061s & 0.073s & 0/3 \\
LLM-evolved & \texttt{condat\_vu} $\tau{=}0.05,\ \sigma{=}1.0,\ \rho{=}1.9$ & \textbf{0.036s} & \textbf{0.073s} & 0/3 \\
\bottomrule
\end{tabular}
\caption{TV-on-OCT cell. Both evolved policies beat the certified textbook
default by $\approx 8.5\times$ on held-out clinical images with zero failures.
The LLM champion's certificate: $\tau L_f/2 + \tau\sigma\ell^2 = 0.425 < 1$,
averagedness $\theta = 0.511$, $\rho\,\theta = 0.97 < 1$.}
\label{tab:tv-results}
\end{table}

The evolved policy is $8.5\times$ faster than the textbook default on held-out
images (Table~\ref{tab:tv-results}), and it is not a standard tuning: it
shrinks the primal step $10\times$, grows the dual step $5.4\times$, retreats
from the stepsize-condition boundary ($0.9925 \to 0.425$), and spends the freed
slack on near-maximal certified over-relaxation ($\rho\,\theta = 0.97$;
Remark~\ref{rem:relaxation-window}). Invalid proposals were rejected
symbolically without consuming solver budget, and no uncertified configuration
was executed. The LLM-guided and random runs tie on held-out quality, and
because they shared a proposal stream we make no sample-efficiency claim from
this cell; Section~\ref{sec:exp:llm} compares against LLM-only evolution
instead. This cell measures evolution against the textbook defaults of the
same scheme, not against a completed field: no external TV solver is run here,
and anisotropic TV-L2 is exactly solvable by parametric max-flow, so the
$8.5\times$ is an internal ablation and not a claim about the state of the art
for this problem.

\paragraph{Operators: certified synthesis.}
Table~\ref{tab:operators} compares implementations of the projection onto a box
intersected with a hyperplane, the operator that relocation needs on the
budget-constrained families. The best synthesized operator is $10.6\times$
faster than the bisection seed the search started from and $3.16\times$ faster
than the textbook breakpoint search on the RTX~4090. A black-box LLM search
without the certificate found an operator that is faster still per call, but
it fails 16 of the 17 contract clauses (maximum error $9.97$) and, used inside
a solver, reaches the target on only $5$ of $9$ unseen instances; the
certified operators reach it on $9$ of $9$, at exactly the iteration count of
the reference projection, on both GPUs.

Per-call speed is not end-to-end speed. With the formulation fixed and only the
projection swapped, on knapsack\_ridge at $n = 12{,}000$ the synthesized
operator's cold first solve is $1.16\times$ to $1.29\times$ faster than with a
hand-written 12-step Newton projection across tolerances $10^{-3}$ to
$10^{-6}$, and $1.27\times$ to $1.56\times$ faster than the textbook breakpoint
search on a warm solve. Most of the speedup over generic solvers comes from the
formulation: relocation with the hand-written Newton projection is already
$3.7\times$ faster than the best engine at $10^{-4}$. Operators evolved in a
later agent session match the earlier synthesized ones to within about $5\%$.

\begin{table}[t]
\centering
\small
\resizebox{\linewidth}{!}{%
\begin{tabular}{lrrcc}
\toprule
projection onto $\{lo \le x \le hi,\ a^{\top}x = b\}$ & RTX 4090 ($\mu$s) & H100 ($\mu$s) & certified & unseen instances \\
\midrule
bisection, 60 steps (search seed) & 395.67 & 439.09 & yes & -- \\
Newton, 12 steps (hand-written) & 165.90 & 113.42 & yes & -- \\
breakpoint search \citep{Kiwiel2008Breakpoint} & 117.88 & 69.71 & yes & -- \\
synthesized, round 1 & 47.54 & 59.44 & yes & 9/9 \\
synthesized, round 2 & \textbf{37.35} & \textbf{33.15} & yes & 9/9 \\
black-box LLM search, no certificate & 12.04 & 8.24 & no (1/17) & 5/9 \\
\bottomrule
\end{tabular}%
}
\caption{Per-call cost of the box-hyperplane projection, from per-device
calibration in the gene pool. ``Unseen instances'' counts held-out instances on
which a solver using the operator reaches the target.}
\label{tab:operators}
\end{table}

\paragraph{Formulations.} The formulation gene moves iteration counts far more
than any parameter: across the formulations of one capacitated routing LP the
iteration count spans $1{,}233\times$, and two formulations whose operator norms
differ by $0.056\%$ differ by $154.7\times$ in iterations. Because the
per-iteration cost also changes, the router decides from measured time, not
from iterations (Section~\ref{sec:limitations}).

\subsection{Wins against structure-specific solvers}

\begin{table}[t]
\centering
\small
\resizebox{\linewidth}{!}{%
\begin{tabular}{llrlrrrr}
\toprule
 & & ours, cold & best engine & \multicolumn{3}{c}{speedup, cold first solve} & warm \\
\cmidrule(lr){5-7}
instance (held-out) & route chosen & $10^{-4}$ (s) & $10^{-4}$ (s) & $10^{-3}$ & $10^{-4}$ & $10^{-6}$ & $10^{-4}$ \\
\midrule
knapsack\_ridge, $n = 12{,}000$ & relocation + synthesized op. & 1.04 & Clarabel 4.69 & $4.08\times$ & $4.49\times$ & $5.00\times$ & $9.17\times$ \\
knapsack\_ridge, $n = 48{,}000$ & relocation + synthesized op. & 8.10 & Clarabel 204.3 & $34.98\times$ & $25.22\times$ & $20.57\times$ & $72.79\times$ \\
robust LP, $n = 5{,}000$ & GPU first-order, standard form & 1.77 & Clarabel 15.71 & $12.64\times$ & $8.86\times$ & $5.63\times$ & $11.41\times$ \\
SVM dual, $n = 20{,}000$ & relocation + synthesized op. & 1.20 & Clarabel 1.24 & $0.84\times$ & $1.03\times$ & $1.11\times$ & $1.97\times$ \\
budget band, $n = 80{,}000$ & relocation + synthesized op. & 1.33 & PIQP 0.32 & $0.24\times$ & $0.24\times$ & $0.25\times$ & $0.42\times$ \\
group lasso, $n = 10{,}000$ & cone relocation & 1.02 & celer 0.57 & $0.57\times$ & $0.56\times$ & $0.55\times$ & $0.88\times$ \\
\bottomrule
\end{tabular}%
}
\caption{The routes the system chose on held-out instances, timed head to head
on an RTX~4090. Ours is a fresh process with detection, build, compilation and
solve in the clock and no persistent compilation cache (cold), or a second solve
in the same process (warm); CUDA context creation, about $0.22$\,s, is outside
the clock. The best engine is the fastest accepted result over every generic
solver and every recognized specialist, each run in a fresh process with a
600\,s limit. Paired geometric means over five repetitions;
intervals are in the experiment log, Section~\ref{LOG-app:field}.}
\label{tab:specialists}
\end{table}

Table~\ref{tab:specialists} times the routes the system chose on held-out
instances. The budget-constrained least-squares family has no packaged solver;
its natural specialist is projected gradient with an exact projection onto the
box and budget constraint, and the system builds that method itself. A cold
first solve is $4.1\times$ to $5.0\times$ faster than every available engine at
$n = 12{,}000$ and $20.6\times$ to $35.0\times$ faster at $n = 48{,}000$, where
Clarabel needs $204$ to $245$\,s and PIQP and ECOS do not finish in $600$\,s. The
margin grows with size: on the same family at tolerance $10^{-4}$ it is
$4.58\times$ at $n = 12{,}000$, $11.37\times$ at $24{,}000$ and $43.09\times$ at
$48{,}000$, and at $n = 100{,}000$ no engine returns an answer (Clarabel exceeds
a 24\,GB memory limit and PIQP returns nothing in $600$\,s) while ours takes
$25.6$\,s. On the robust LP the evolved route is $12.6\times$ faster at
$10^{-3}$ and $5.6\times$ at $10^{-6}$. On three of the six instances the route
the system chose loses when timed head to head: the SVM dual is at parity or slower,
PIQP is $4\times$ faster on the budget band, and celer is about $1.8\times$
faster on group lasso. The router's earlier cost predictions for these
instances were too optimistic, which Section~\ref{sec:limitations} discusses.

\paragraph{Lasso paths with dense solutions.} On a 100-value lasso
regularization path, scored per regularization value by the lasso duality gap,
with celer and scikit-learn each run as shipped, the batched
GPU method wins once the solution support approaches the number of rows. At
$n = 10{,}000$ and tolerance $10^{-4}$ the speedup grows with density from
$0.52\times$ to $1.52\times$ to $15.8\times$ ($3.28$\,s against $51.80$\,s), and
at $10^{-6}$ it is $5.07\times$ ($18.27$\,s against $92.68$\,s). At
$n = 50{,}000$ it is $7.34\times$ at $10^{-4}$ ($29.55$\,s against
$216.92$\,s), and at $10^{-6}$ our method certifies all $100$ values in
$177$\,s while the best completed specialist run certifies $72$ in $641$\,s.

\paragraph{Other specialist comparisons.} On the SVM family at
$n_{\mathrm{var}} = 110{,}000$ and tolerance $10^{-6}$, LIBLINEAR
\citep{Fan2008LIBLINEAR} takes $43.0$\,s on the identical instance, against
$10.35$\,s for our method in the first-reported measurement; under the corrected
accounting our margin over Clarabel on this family is $3.78\times$
$[3.67, 3.90]$. On elastic-net logistic regression over the epsilon data
set, the method is $2.28\times$ $[1.73, 3.01]$ faster than the GPU specialist
cuML \citep{NVIDIA2026cuML} on ten of ten penalty settings for a repeated
solve, and $0.06\times$ on a first solve with an empty compilation cache.

\paragraph{Across the held-out set.} Over 28 held-out instances spanning QP,
LP, SOCP, Maros--M\'esz\'aros, Netlib and CBLIB, the routed system solves $27$
within a 20\,s budget, against $26$ for the best-per-instance oracle over the
completed field and $25$ for the best single engine (Clarabel). On the $25$
instances both solve, it is faster on $13$ and the oracle on $12$, with a
geometric-mean speedup of $1.29\times$ on repeat visits and $0.035\times$ on
first visits, when exploration dominates.

\subsection{Against general-purpose solvers at scale}

\begin{table}[t]
\centering
\small
\begin{tabular}{lrrl}
\toprule
family & $n_{\mathrm{var}}$ & speedup over Clarabel & $95\%$ interval \\
\midrule
SVM dual & 550{,}000 & $7.91\times$ & $[6.34, 9.88]$ \\
SVM & 110{,}000 & $3.78\times$ & $[3.67, 3.90]$ \\
portfolio & 500{,}707 & $2.37\times$ & $[1.78, 3.16]$ \\
\bottomrule
\end{tabular}
\caption{Speedup over Clarabel at tolerance $10^{-6}$, three seeds, every
repetition paying its own setup, under the corrected accounting of the
experiment log, Section~\ref{LOG-app:accounting}. Families appear here only where
Clarabel is the strongest available engine. Huber, lasso and the minimum
enclosing ball are excluded because their structure-specific solvers are faster
than our method (Section~\ref{sec:limitations}), and the control family is
excluded because OSQP, already in our own field, is faster than Clarabel at that
size.}
\label{tab:general}
\end{table}

Against a general-purpose interior-point solver, the evolved formulations win
on every family in Table~\ref{tab:general}, and every interval excludes $1$.
The wins are size dependent: the SVM dual crosses over between
$n_{\mathrm{var}} = 38{,}500$ ($0.79\times$) and $55{,}000$ ($1.45\times$), and
portfolio between $100{,}316$ ($0.92\times$) and $250{,}500$ ($1.57\times$).

\subsection{Against LLM-only evolution}



The LLM-only arm evolves an operator with the same LLM and the same search
engine, but without the certificate, the gene pool or the router, and uses it in a
fixed pipeline. It solves $18$ of $28$ held-out instances against $27$ for the
full system, and its geometric-mean repeat solve is $23\times$ slower
(Table~\ref{tab:llm}). Its first visit is slightly faster because it has a
single candidate to run. The faster uncertified operator passes only 1 of 17
contract clauses and reaches the target on 5 of 9 unseen instances, against 9
of 9 for the contract-checked operators (Table~\ref{tab:operators}). LLM token
counts were not logged per arm, so the two arms are matched on search
iterations, not on tokens.

\subsection{Choosing CPU or GPU}

The router's device choices follow instance structure and size. It sends
sparse lasso to a CPU coordinate-descent specialist (celer, a routed cost of
$11.9$\,ms at $n = 5{,}000$), small Maros--M\'esz\'aros and Netlib instances to
a CPU interior-point method (about $1$\,ms), and large structured QPs and robust
LPs to evolved GPU methods (Table~\ref{tab:specialists}). On MIPLIB LP
relaxations, $20$ of the $22$ instances solved within budget are routed to CPU
interior-point engines and $2$ to our first-order formulations. The
payoff of a formulation also depends on the device: on the capacitated routing
LP the simplex formulation's advantage is $5.6\times$ on the GPU and $45.0\times$
on the CPU, and one SVM instance loses at $0.41\times$ on the H100 and wins at
$1.22\times$ on the RTX 4090 because the dominant work is single-threaded host
code and the H100 node has the weaker single thread.

\begin{table}[t]
\centering
\small
\begin{tabular}{lccc}
\toprule
device, precision & tolerance $10^{-4}$ & $10^{-6}$ & $10^{-8}$ \\
\midrule
H100, float64 & 23.01\,s & \textbf{45.60\,s} & \textbf{76.88\,s} \\
H100, float32 & \textbf{9.69\,s} & never & never \\
RTX 4090, float64 & 35.64\,s & 71.77\,s & 121.83\,s \\
RTX 4090, float32 & 16.52\,s & never & never \\
\bottomrule
\end{tabular}
\caption{The same batched lasso cross-validation workload ($n = 50{,}000$, 5
folds, 100 regularization values) on two devices and two precisions. float32
stalls at a relative gap of about $2\times10^{-5}$, so the winning pair changes
with the target.}
\label{tab:device}
\end{table}

Table~\ref{tab:device} shows that the winning pair of device and precision
changes with the accuracy target. At $n = 10{,}000$ and $10^{-4}$ the consumer
RTX 4090 in float32 is the fastest of the four ($1.63$\,s). The H100's float64
advantage over the RTX 4090 is $1.6\times$, far below the ratio of their float64
throughput and below their $3.3\times$ memory-bandwidth ratio. On this
cross-validation workload the CPU specialists are faster still, which is the
routing decision the system should and does make.

The kernel a device should run is also a per-device decision. Sweeping 21
temporally fused CUDA tiles against XLA's compiled loop on an RTX 4090 and an
H100, the best tile differs at five of six problem sizes, the fused kernel is
$1.15\times$ to $1.66\times$ faster than XLA on the 4090 and $1.95\times$ to
$2.74\times$ faster on the H100, and carrying one device's choice to the other
costs $1.06\times$ to $1.63\times$. A selector that reads the CUDA device
properties reproduces each machine's choice from one table with no recompile,
and the kernel is bitwise identical to the unfused iteration at relaxation $1$
on both devices (experiment log, Section~\ref{LOG-app:log:fusion}).

\paragraph{The device is a priced candidate, and the crossover is measurable.}
A candidate is a triple (formulation, operator, device): the same formulation on
this machine's GPU and on its CPU are two candidates, sharing one measured
iteration count and priced from each device's own calibration, with a device
this process has no backend for declined rather than guessed at. Measuring our
method on both devices for all 28 held-out instances, four repetitions, the
iteration count is device independent as the cost model assumes (median ratio
$1.0000$), and the two machines disagree about which device to use: at
tolerance $10^{-3}$ the GPU is faster on 18 instances and the CPU on 2 on the
RTX 4090, against 14 and 6 on the H100.

What predicts the crossover is not the instance size. Ranked by $\mathrm{nnz}(A)$,
13 of 27 instances fall on the wrong side of any threshold, because relocation
removes rows from the operator the iteration actually applies. Ranked by the
work one iteration does, $\mathrm{nnz}(L) + n$, the CPU has the cheaper
iteration up to about $14{,}400$ and the GPU from about $22{,}400$, and on the
RTX 4090 no instance falls on the wrong side of that band. Below it the GPU is
launch bound by up to $600\times$ per iteration; above it the GPU wins by up to
$142\times$. Setup does not discriminate between the devices, which is why
several small instances tie on total time despite a large per-iteration gap.

With device selection enabled the router picks the faster device on 14 of the 15
decidable instances, and on 9 of 9 when it chose the same formulation and
operator the ground truth timed, so only the device differs; its one mistake
costs $1.37\times$. The honest accounting is that this buys little on this
workload: repeat-visit time improves by $2\%$ while first-visit wall time rises
$1.70\times$ from the extra exploration, and 22 of the 28 instances already
route to a CPU engine by recognition rather than by the device axis. The cost
model is biased low on both devices, predicting $0.75\times$ of measured time on
the GPU and $0.63\times$ on the CPU, but the bias is common enough that it still
orders the two devices correctly on 15 of 15 instances, which is what the
decision requires.
